% native creation stories: ancient

\subsection{Greek}

In the beginning there was only \textbf{Chaos}, a yawning gap with nothing
in it. No maker stood behind it. It simply came first. Out of Chaos, on their
own and begotten by no one, the first powers came into being. \textbf{Gaia},
the broad earth, appeared, the seat of all things. \textbf{Tartaros}, the misty
pit, opened far below. And \textbf{Eros}, desire, appeared too, the force of
joining that lets everything else come together and breed.

From Gaia, alone and without a mate, came \textbf{Ouranos}, the starry sky,
stretched over her as her equal and her cover. Gaia and Ouranos were the first
parents. They brought forth the \textbf{Titans}. Ouranos hated his offspring and
pressed them back down into the earth so none could be born. Gaia, in pain, made
a sickle of adamant and asked her children to act. The youngest and boldest,
\textbf{Kronos}, ambushed his father and cut him down, taking his place as the
new ruler.

Kronos feared the same fate. A warning said one of his own children would unseat
him, so he swallowed each child born to him and \textbf{Rhea} as it came. Rhea
hid the last one, \textbf{Zeus}, away on Crete and gave Kronos a stone to swallow
instead. Zeus grew in secret, returned, and forced his father to bring up the
swallowed children. A long war between the generations, the \textbf{Titanomachy},
followed. The young gods won. They cast the Titans down into Tartaros.

Zeus then ordered the world. He and his two brothers divided the realms by lot.
Zeus took the sky, \textbf{Poseidon} took the sea, and \textbf{Hades} took the
world of the dead below. The earth they held in common. The cosmos settled under
the rule of Zeus.

\subsection{Roman}

In the beginning there was a formless \textbf{Chaos}, a rough and unordered heap
in which the elements warred against one another. Earth, water, air, and fire
were jumbled together, each opposing the rest. There was no shape and no settled
place for anything.

Then a shaping power, in \textbf{Ovid}, whichever of the gods it was, ended the
conflict. It separated \textbf{Tellus} the earth from \textbf{Caelus} the sky,
the sea from the land, and the heavy from the light. It rounded the earth, set
the seas in their basins, lifted the air, and placed the stars. Last of all came
man, made to look up at the sky. The warring heap became an ordered world with
its realms laid out in order. \textbf{Janus}, the two-faced god of doorways and
beginnings, stood at the very start of things.

After this came the ages of mortals, a long decline. First a Golden Age of ease
and peace under \textbf{Saturn}, with no toil and no law. Then a Silver Age, then
a Bronze Age, and at last a hard Iron Age of labor, greed, and war. The gold of
the first age gave way step by step to the iron of the present one.

Behind this ordering lay the borrowed genealogy. \textbf{Saturn} had unseated
\textbf{Caelus} and swallowed his own children by \textbf{Ops}, until
\textbf{Jupiter}, hidden away, grew and overthrew him. Jupiter took the sky,
\textbf{Neptune} the sea, and \textbf{Pluto} the underworld. The same arc as the
Greek runs underneath the Roman world.

\subsection{Norse}

In the beginning there is no earth and no sky. There is only \textbf{Ginnungagap},
a vast yawning emptiness. On one side lies \textbf{Niflheim}, a realm of frozen,
unformed matter. On the other lies \textbf{Muspelheim}, a realm of fire and heat.
Cold flows in from one edge and heat flows in from the other.

Where the cold and the heat meet in the middle of Ginnungagap, the first being,
\textbf{Ymir}, takes shape from the meltwater. He is a single immense frost-giant
holding the raw matter of the world. From his own sweat more giants come. The cow
\textbf{Audhumla} also forms from the thaw, and by licking the salty ice she
brings to light \textbf{Buri}, the ancestor of the gods.

From Buri descend three shaping powers, \textbf{Odin}, \textbf{Vili}, and
\textbf{Ve}. The three slay Ymir. They carry his body into the middle of
Ginnungagap and build the world out of it. The flesh becomes the ground. The
blood becomes the seas. The bones become the mountains. The skull becomes the
sky. The brains become the clouds.

The shapers then find two pieces of driftwood washed up at the edge of the world.
They give them breath, sense, and warmth, and these become \textbf{Ask} and
\textbf{Embla}, the first humans. The whole structure hangs on \textbf{Yggdrasil},
the great ash-tree whose trunk is the axis and whose roots reach the deep wells.
The order is bound by the law of fate, and that law foretells \textbf{Ragnarok},
a great ending in which the powers fall in battle and fire consumes the world.
After the ending the ground rises green and new from the sea, and life returns.

\subsection{Egyptian}

In the beginning there is only \textbf{Nun}, an endless dark ocean. It is not
nothing. It is undifferentiated potential, formless and motionless, holding
everything that could be.

From the waters the first mound of dry land, the \textbf{Benben}, rises. On the
mound \textbf{Atum} appears, the self-made one. He is complete in himself and
needs no partner. By his own substance he brings forth the first pair,
\textbf{Shu}, the principle of dry expansive air, and \textbf{Tefnut}, the
principle of moist receptivity. From that pair come the next pair, \textbf{Geb}
the ground beneath and \textbf{Nut} the boundary above.

At first Geb and Nut are locked together. Shu lifts Nut off Geb and holds them
apart, opening the room in which all things can exist. From Geb and Nut come the
ruling powers, the \textbf{Ennead}, the generation that governs the world,
including \textbf{Osiris} the just king, \textbf{Isis} the worker of magic, and
\textbf{Set}, the force of disruption that opposes the order.

What Atum establishes is \textbf{Maat}, order itself, the standard of truth and
balance. Against it presses \textbf{Isfet}, disorder, which never fully departs
because the waters of Nun never fully go away. The tradition as a whole is the
continual work of holding Maat in place against Isfet. The clearest daily image
of this is the journey of \textbf{Ra} across the sky and through the \textbf{Duat}
below, attacked each night by the serpent \textbf{Apep} and reborn each dawn. In
the Memphite telling \textbf{Ptah} stands prior even to Atum, making all things
by the thought of his heart and the word of his tongue.

\subsection{Mesopotamian}

In the beginning there was only water. Two waters mingled together.
\textbf{Apsu} was fresh, calm and sweet. \textbf{Tiamat} was salt, vast and
churning. Where the two waters met and ran into each other, the first gods were
born.

The young gods were loud. Their noise broke the rest of the deep. Apsu resolved
to kill them and have peace again. But \textbf{Ea}, the cleverest of the young
gods, learned the plot first. He put Apsu to sleep and slew him, and he built his
own house upon the slain body.

Then Tiamat rose in grief and rage. She bred an army of monsters and set
\textbf{Kingu} over them as her champion, and she gave him the \textbf{Tablet of
Destinies}, the warrant of supreme rule. The gods were afraid. None would face
her. At last the young storm-god \textbf{Marduk} stepped forward. He asked for
kingship over all the gods as his price, and they granted it.

Marduk went out against Tiamat with the winds and a net and a bow. He drove the
winds into her open mouth so she could not close it, and he loosed his arrow and
split her. Then he cut her body in two. From the upper half he made the sky, and
from the lower half the earth. From her eyes he set the Tigris and the Euphrates
flowing. He fixed the stars and the moon and the sun in their courses and laid
down the calendar.

The gods were still weary, for someone had to do the digging and the work. So
Marduk took the blood of Kingu, mixed it with clay, and made humanity. Humans
were made to carry the labor of the gods and to feed them with offerings. Then
the \textbf{Anunnaki} built the city of Babylon for their king and proclaimed his
fifty names.

\subsection{Chinese}

In the beginning there was only \textbf{Hundun}, a formless chaos shaped like an
egg. All things were mingled together inside it, \textbf{yin} and \textbf{yang},
the light and the heavy, the clear and the murky, none of them yet apart. Within
the egg \textbf{Pangu}, the first being, slept and grew for a vast age.

When Pangu woke, he broke the egg open. The light and clear yang floated up and
became the sky. The heavy and murky yin sank down and became the earth. Pangu
stood between them and held them apart. Each day the sky rose higher, the earth
grew thicker, and Pangu grew taller to match, pushing the two ever further apart
so they could not fall back together. This went on for another vast age.

When the work was done, Pangu lay down and died. His body became the world. His
breath became the wind and clouds. His voice became the thunder. His two eyes
became the sun and the moon. His limbs became the mountains, his blood the
rivers, his flesh the soil, and his hair the stars and the plants.

The world was made but empty of people. \textbf{Nuwa}, the serpent-bodied mother,
walked the new earth. She took yellow clay and modeled the first humans by hand
in her own image. When she tired of shaping each one by hand, she dipped a cord
into the mud and flung it, and every drop that fell became a person. So the world
was filled. When the pillar of heaven later broke, Nuwa melted five-colored
stones to patch the sky and propped the four corners on the legs of a great
turtle.
